\begin{figure}[h]
\centering
\begin{tikzpicture}[
    scale=1.3,
    >=Stealth, % 使用箭头样式
    font=\small
]

    % --- 定义坐标 ---
    % 定义原点 Q (控制结果)
    \coordinate (Q) at (0,0);

    % 定义流形直线的角度 (例如 15度)
    \def\ang{15}

    % 定义流形上的点 (缩短了长度)
    % R 点的位置 (理想态)
    \coordinate (R) at ($(Q) + (180+\ang:2.0)$);

    % 流形起点 (稍微延伸出 R 点一点点)
    \coordinate (M_start) at ($(Q) + (180+\ang:2.5)$);
    % 流形终点 (稍微延伸出 Q 点一点点)
    \coordinate (M_end)   at ($(Q) + (\ang:1.5)$);

    % 定义 P (现状) 在垂线上 (视觉垂直)
    \coordinate (P) at ($(Q) + (90+\ang:2.2)$);

    % --- 绘制图形 ---

    % 1. 绘制流形 M_target (蓝色粗线，缩短版)
    \draw[blue, very thick] (M_start) -- (M_end) node[right, black] {$\mathcal{M}_{target}$ (目的流形)};

    % 2. 绘制 m-投影线 P -> Q (虚线, 带箭头)
    \draw[dashed, thick, ->] (P) -- (Q);

    % 3. 绘制 P -> R 的连接线 (细虚线)
    \draw[dashed, thin] (P) -- (R);

    % 4. 绘制直角标记
    \draw[thick] ($(Q)!0.15cm!(P)$) --
                 ($(Q)!0.15cm!(P)!0.15cm!90:(P)$) --
                 ($(Q)!0.15cm!(M_end)$);

    % --- 添加点和标签 ---

    % 点 P
    \fill (P) circle (1.5pt) node[above] {$P$ (现状)};

    % 点 Q
    \fill (Q) circle (1.5pt) node[below=3pt] {$Q$ (控制结果)};

    % 点 R
    \fill (R) circle (1.5pt) node[below=3pt] {$R$ (理想态)};

    % --- 添加说明性文字 ---
    \node[right, align=left] at ($(Q)!0.5!(P)$) {
        $\nabla^{(m)} \perp \nabla^{(e)}$ \\
        \footnotesize{m-投影}
    };
\end{tikzpicture}
\caption{投影关系的几何示意；散度不等式与等式分别需要本章所述的凸性和正交条件}
\end{figure}
